\clearpage
\Article{358.\ 359}
\textsc{Satz.} \emph{Es sei $P$ eine Primfunktion und $x^\nu$ die niedrigste Potenz von $x$, die bei der Division durch $P$ den Rest $1$ ergibt. Ferner sei $P\equiv(P,\rho^n)$. Dann ist $n$ einer Potenz der Zahl $p$ nach dem Modul $\nu$ kongruent.}

\emph{Beweis.} Wir haben oben gezeigt, dass, wenn $P$
\[
=x^m+Ax^{m-1}+Bx^{m-2}+\text{usw.}
\]
ist, der Ausdruck
\[
\begin{aligned}
z^m+Az^{m-1}+Bz^{m-2}+\text{usw.}\qquad\qquad\qquad\qquad&\\
{}-(z-x)(z-x^p)\dotstwo(z-x^{p^{m-1}})&
\end{aligned}
\]
durch $P$ teilbar ist. Auf dieselbe Weise würde folgen, dass
\[
\begin{aligned}
z^m+Az^{m-1}+Bz^{m-3}+\text{usw.}\qquad\qquad\qquad\qquad&\\
{}-(z-x^n)(z-x^{np})\dotstwo(z-x^{np^{m-1}})&
\end{aligned}
\]
durch $P$ teilbar ist. Da diese Faktoren aber untereinander teilerfremd sind, müssen die einzelnen Faktoren notwendig den entsprechenden einzelnen Faktoren nach $P,p$ kongruent sein. Daher muss $z-x^n$ $\equiv z-x^{p^\kappa}$ sein, d. h. $p^\kappa\equiv n(\mathrm{mod.}\,\nu)$. Q. E. D.\begingroup\interfootnotelinepenalty=10000\relax\renewcommand{\thefootnote}{*)}\footnote{Ist $(P,\rho^a)\equiv(P,\rho^b)(\mathrm{mod.}\,p)$, so ist $a\equiv p^\kappa b(\mathrm{mod.}\,\nu)$.

\emph{Beweis.} Es sei $z^m+Az^{m-1}+Bz^{m-2}+\dotstwo=\Pi$. Dann ist $(\Pi,\rho^a)\equiv(\Pi,\rho^b)(\mathrm{mod.}\,P)$; es gilt aber
\[
\begin{aligned}
(\Pi,\rho^a)&\equiv(z-x^a)(z-x^{ap})(z-x^{app})\dotstwo(z-x^{ap^{m-1}}),\\
(\Pi,\rho^b)&\equiv(z-x^b)(z-x^{bp})(z-x^{bpp})\dotstwo(z-x^{bpm-1})
\end{aligned}
\]
woraus die Behauptung folgt.

Das Produkt aus $\Pi,(\Pi,\rho^2),(\Pi,\rho^3)$ usw. $(\Pi,\rho^\nu)$ ist $\equiv(z^\nu-1)^m(\mathrm{mod.}\,P)$, denn es ist
\[
\begin{aligned}
(z-x)(z-x^2)(z-x^3)\ldots(z-x^\nu)&\\
{}\equiv(z-x^p)(z-x^{2p})(z-x^{3p})\dotstwo(z-x^{\nu p})&\equiv\text{usw.}\equiv z^\nu-1
\end{aligned}
\]
In der Folge $P,(P,\rho^2),(P,\rho^3)$ usw. $\dotsfour(P,\rho^\nu)$ kommen alle Primteiler der Funktion $x^\nu-1$ vor, und zwar jeder $m$ Mal. Daraus folgt, dass das Produkt aus allen $\equiv(x^\nu-1)^m$ ist.}\endgroup

\par\medskip\Needspace{5\baselineskip}
\begin{center}\emph{Die Primteiler der Funktion $x^\nu-1$ nach einem Primzahlmodul bestimmen.}\end{center}
\Article{360}
Ist $\nu$ durch den Modul $p$ oder eine seiner Potenzen teilbar, so sei $\nu=p^k\lambda$. Dann ist
\[
x^\nu-1\equiv(x^\lambda-1)^{p^k}(\mathrm{mod.}\,p)
\]
Hieraus ergibt sich, dass nur der Fall betrachtet werden muss, in dem $\nu$ nicht durch $p$ teilbar ist.

\SourceLeaf{240}
Ist $p^m\equiv1(\mathrm{mod.}\,\nu)$, und zwar $m$ so klein wie möglich, so ist offenbar $x^{p^m-1}-1$ durch $x^\nu-1$ teilbar. Daher kann $x^\nu-1$ nur solche Teiler haben wie $x^{p^m-1}-1$. Dieser Ausdruck hat aber Primteiler vom Grad $m$ und andere, deren Grad ein Teiler der Zahl $m$ ist. Solche Teiler hat daher auch $x^\nu-1$. Wie viele von jeder Art vorhanden sind, zeigen wir an einem Beispiel, aus dem sich das allgemeine Gesetz leicht ableiten lässt.

Es sei $\nu=63$ und $p=13$. Dann ist $m=6$. Daher hat $x^{63}-1$ nach dem Modul $13$ Primfaktoren der Grade sechs, drei, zwei und eins. Offenbar ist das Produkt aus den Faktoren ersten Grades der gemeinsame Teiler höchsten Grades der Funktionen $x^{63}-1$ und $x^{12}-1$, d. h. $x^3-1$. Es gibt also drei Primfaktoren ersten Grades. Das Produkt aus allen Primfaktoren zweiten und ersten Grades ist der gemeinsame Teiler der Funktionen $x^{63}-1$ und $x^{168}-1$, d. h. $x^{21}-1$. Daher gibt es $\frac{21-3}2$ oder $9$ Faktoren zweiten Grades. Das Produkt aus den Primfaktoren dritten und ersten Grades ist der gemeinsame Teiler der Funktionen $x^{63}-1$ und $x^{2196}-1$, d. h. $x^9-1$. Daher gibt es $\frac{9-3}3$, d. h. $2$ Teiler dritten Grades. Die übrigen haben schließlich den Grad sechs; ihre Anzahl ist daher $=\frac{63-6-18-3}6$, d. h. $6$.

\Needspace{12\baselineskip}
Durch genaue Betrachtung dieses Sachverhalts lässt sich leicht die folgende allgemeine Regel herleiten:

\emph{Es sei $\delta$ ein Teiler von $m$, und es seien alle Teiler der Zahl $\delta$, die kleiner als $\delta$ sind, $\delta',\delta'',\delta'''$ usw. Die größten gemeinsamen Teiler von $\nu$ mit $p^\delta-1$, $p^{\delta'}-1$, $p^{\delta''}-1$ usw. seien der Reihe nach $\mu,\mu',\mu''$ usw. Es sei $\frac\mu{\mu'};\frac\mu{\mu''},\frac\mu{\mu'''}$ usw. $=\lambda',\lambda'',\lambda'''$ usw. Dann hat $x^\nu-1$ $\frac1\delta$ Mal so viele Primteiler vom Grad $\delta$, wie es unterhalb von $\mu$ Zahlen gibt, die durch keine der Zahlen $\lambda',\lambda'',\lambda'''$ usw. teilbar sind.}

\Needspace{10\baselineskip}
\Article{361}
\textsc{Satz.} \emph{Ist eine Funktion $X$ der Unbestimmten $x$ durch eine andere Funktion $\xi$ teilbar und geht $X$, wenn man $x$ durch $x^k$ ersetzt, in $X'$ über, so ist $X'$ durch $(\xi,\rho^{\frac1k})$ teilbar.}

\emph{Beweis.} Es sei $X\equiv\xi\upsilon$, und $\xi,\upsilon$ mögen in $\xi',\upsilon'$ übergehen, wenn man $x$ durch $x^k$ ersetzt. Offenbar ist dann $X'\equiv\xi'\upsilon'$. Doch ist $\xi'$ durch $(\xi,\rho^{\frac1k})$ teilbar. Also gilt dies auch für $X'$. Q. E. D.
